% Fragmento público generado desde propietarios íntegros.
% Residencia: 52.3.1.
% Fuente: ../incoming/radion_monografia_20260827/manuscrito/sections/04_cincuenta_critica.tex:1-105
\noindent The critical coordinate links the Cayley chart, the APP record, and the dodecaphasic phase through typed maps.\par\medskip

\subsubsection{Two exact realizations of the angular coordinate \(50^\circ\)}

\paragraph{Second phase of the dodecaphasic temporal system}

By \eqref{eq:rueda-dodecafasica}, the sequence of centros begins
\[
20^\circ,\ 50^\circ,\ 80^\circ,\ 110^\circ,\ldots.
\]
Solve
\[
30m-10=50
\]
gives only \(m=2\). Therefore, \(50^\circ\) is not a rounding of \(A\),
a half-turn, or the centre of a circle: it is a typed phase of the
dodecaphase calendar. The fraction of a turn is
\[
\frac{50}{360}=\frac5{36},
\]
exactly the first summand of \(S_E/T_+\).

\paragraph{Cayley Critical Coordinate Matching}

The functional involution
\[
\rho\longmapsto1-\bar\rho
\]
fixed
\[
\operatorname{Fix}(\rho\mapsto1-\bar\rho)
=\{\rho:\Re\rho=1/2\}.
\]
The number \(1/2\) is the self-dual real coordinate of the functional interval \(0\leq\Re s\leq1\). In the longitudinal sense, it is not \enquote{half of a line}, and it does not contain the spectral height \(\gamma\). For
\[
\rho_\gamma=\frac12+i\gamma,
\]
The Cayley transformation can be written
\begin{equation}
\tau_\gamma
=\frac{\rho_\gamma-1}{\rho_\gamma}
=\frac{\gamma+i/2}{\gamma-i/2}
\in S^1,
\qquad
\gamma=\frac12\cot\frac{\theta_\gamma}{2}.
\label{eq:cayley-critica}
\end{equation}
Thus, \(1/2\) fixes the seam, \(\gamma\) determines the angular character, and
\(\Gamma_9\) preserves the depth of its iterations.

\paragraph{Affine injection \(j_{100}\) of the APP record into the angular chart}

Cubic completion arranges a chamber of ten marks per coordinate. One
face contains \(10^2=100\) positions. We declare the affine chart
\begin{equation}
j_{100}:[0,1]\longrightarrow[0,100^\circ],
\qquad j_{100}(\sigma)=100\sigma^\circ.
\label{eq:j100}
\end{equation}
Then
\[
j_{100}\left(\frac12\right)=50^\circ.
\]
Requiring \(j_{100}(1/2)=\theta_m\) simultaneously yields
\[
100\cdot\frac12=30m-10
\quad\Longleftrightarrow\quad m=2.
\]

\begin{theorem}[Critical Equalizer--Dodecaphasic]
In the hundred-mark APP chart, the fixed coordinate of the critical involution
is published exactly in the second phase of the wheel:
\begin{equation}
\boxed{j_{100}(1/2)=\theta_2=50^\circ.}
\label{eq:igualador-50}
\end{equation}
The phase \(m=2\) is unique.
\end{theorem}

The theorem is exact once \eqref{eq:j100} has been declared. Its status is
\status{new formalization}: the critical midpoint, the \(10^3\) completion, the hundred-mark face, and \(\theta_2\) had already been constructed; choosing that face as the affine critical chart brings their domains together for the first time.

\paragraph{Compatibility of the critical value with the Cayley chart}

The Cayley transform does not send \(1/2\) to \(50^\circ\). If
\(\gamma=0\), \eqref{eq:cayley-critica} gives
\[
\tau_0=-1=e^{i\pi},
\]
that is, \(180^\circ\). The \(50^\circ\) belong to the APP chart of one hundred
marks and to the dodecaphasic calendar; the Cayley angle depends on \(\gamma\).
This distinction prevents all the zeros of the zeta function from being collapsed into a single
phase.

\begin{falsifier}
The identification \(1/2\leftrightarrow50^\circ\) fails if the
chart \eqref{eq:j100} is removed. The scalar equality \(100(1/2)=50\) is not sufficient
by itself to identify two objects. The theorem depends on the distinguished
APP face and on the dodecaphase wheel.
\end{falsifier}

